\section{Experiments}

\noindent{\textbf{Datasets.}}
We train \OURMODEL\ on the large-scale RealEstate10K (RE10K)~\cite{zhou2018re10k} and ACID~\cite{liu2021acid} datasets, following prior works for train/test splits.
For two-view evaluation, we adopt the same test split as prior feed-forward methods~\cite{charatan2024pixelsplat, chen2024mvsplat, xu2025depthsplat, li2025vicasplat, ye2024nopo}. 
For multi-view evaluation, we use the scene categorization provided by NoPoSplat~\cite{ye2024nopo} and select input pairs with small overlap. Additional input views are sampled between the selected pair without duplication to match the target number of input views (8, 16, and 24 views).

\noindent{\textbf{Implementation Details.}} We use three levels of multi-resolution feature maps (\(L=3\)) in all experiments and choose the finest level to match the input image resolution, $(H^L, W^L)\!=\!(H,W)$. We initialize our model from pretrained VGGT weights. We use a learning rate of \(2\times10^{-4}\) for most modules, while freezing the patch embedding weights of the geometry backbone and training the remaining backbone parameters with a smaller learning rate of \(2\times10^{-5}\).
For multi-view training, we train on RE10K. At each iteration, each GPU independently samples the number of input views from $\{2,3,4,6,12,24\}$ and selects context images accordingly. We additionally sample the same number of target novel views for supervision. To keep the total number of images processed per iteration constant, we use a dynamic batch size that is inversely proportional to the number of context images. The multi-view model is trained for 15,000 iterations.
For two-view training, separate models are trained on RE10K and ACID, following the training configuration of NoPoSplat~\cite{ye2024nopo}. Each model is trained for 18,750 iterations with a total batch size of 128. We set the weight of MSE and LPIPS loss as 1 and 0.05 for $\mathcal{L}^\text{render}$. For the final total loss, we fix the weight of $\mathcal{L}^\text{render}$, $\mathcal{L}^\text{score}$, $\mathcal{L}^\text{camera}$, and $\mathcal{L}^\text{scene}$ as $1.0$, $10^{-4}$, $10.0$, and $10^{-2}$ respectively, in all experiments. All experiments are conducted on eight NVIDIA H200 GPUs, and each training run takes approximately 15 hours.

\noindent{\textbf{Baselines and Evaluation Metrics.}}
For multi-view evaluation, we compare \OURMODEL\ with recent feed-forward 3DGS methods under two settings: \textit{uncalibrated} methods that take only images as input without camera poses or intrinsics, and \textit{pose-free} methods that do not require camera poses but assume known intrinsics. 
For the uncalibrated setting, we compare against VicaSplat~\cite{li2025vicasplat} and AnySplat~\cite{jiang2025anysplat}. For fair comparison, we retrain AnySplat under the same training setup as ours, which requires approximately 27 hours on eight NVIDIA H200 GPUs. VicaSplat uses the officially released multi-view pretrained weights on the RE10K dataset. 
For the pose-free setting, we compare against NoPoSplat~\cite{ye2024nopo}, VicaSplat~\cite{li2025vicasplat}, and SPFSplat~\cite{huang2025spfsplat}, using their officially released pretrained models.
For two-view evaluation, we compare with pixelSplat~\cite{charatan2024pixelsplat}, MVSplat~\cite{chen2024mvsplat}, NoPoSplat~\cite{ye2024nopo}, VicaSplat~\cite{li2025vicasplat}, and SPFSplat~\cite{huang2025spfsplat}. We also compare baselines that are not a feed-forward 3D Gaussian splatting method, including pixelNeRF~\cite{yu2021pixelnerf}, DUSt3R~\cite{wang2024dust3r}, MASt3R~\cite{leroy2024mast3r}, and CoPoNeRF~\cite{hong2023coponerf}.
We evaluate the performance of novel view synthesis using three standard metrics: LPIPS, SSIM, and PSNR. We also report the number of final Gaussian primitives to enable joint comparison of rendering quality and Gaussian-count efficiency.
\vspace{-10pt}

\begin{table*}[t]
\centering
\setlength{\tabcolsep}{0.09cm}
\renewcommand{\arraystretch}{0.95}
\caption{\textbf{Novel view synthesis performance on RE10K~\cite{zhou2018re10k} under different numbers of input views.}
We report LPIPS/SSIM/PSNR for 8, 16, and 24 input views.
Best and second-best results are highlighted in \textbf{bold} and \underline{underlined}, respectively. $\tau^+$ and $\tau^-$ denote the high- and low-threshold variants, respectively.}
\label{tab:re10k_multiviews}
\vspace{-6pt}

\resizebox{\textwidth}{!}{
\begin{tabular}{llcccc cccc cccc}
\toprule
& \multirow{2}{*}{\textbf{Method}}
& \multicolumn{4}{c}{\textbf{8 views}}
& \multicolumn{4}{c}{\textbf{16 views}}
& \multicolumn{4}{c}{\textbf{24 views}} \\
\cmidrule(lr){3-6} \cmidrule(lr){7-10} \cmidrule(lr){11-14}
& 
& \#GS$\downarrow$ & LPIPS$\downarrow$ & SSIM$\uparrow$ & PSNR$\uparrow$
& \#GS$\downarrow$ & LPIPS$\downarrow$ & SSIM$\uparrow$ & PSNR$\uparrow$
& \#GS$\downarrow$ & LPIPS$\downarrow$ & SSIM$\uparrow$ & PSNR$\uparrow$ \\
\midrule

\multirow{3}{*}{\shortstack[l]{\emph{Pose-Free}}}
& NoPoSplat
& 524K & 0.213 & 0.756 & 22.31
& 1049K & 0.252 & 0.713 & 21.11
& 1573K & 0.275 & 0.691 & 20.49 \\

& VicaSplat
& 524K & 0.241 & 0.713 & 21.18 
& 1049K & 0.384 & 0.596 & 17.56
& 1573K & 0.443 & 0.546 & 16.13 \\

& SPFSplat
& 524K & 0.142 & 0.849 & 25.66
& 1049K & 0.137 & 0.855 & 25.88
& 1573K & 0.139 & 0.853 & 25.78 \\

\midrule

\multirow{4}{*}{\shortstack[l]{\emph{Uncalibrated}}}
& VicaSplat
& 524K & 0.258 & 0.686 & 20.77
& 1049K & 0.417 & 0.556 & 16.78
& 1573K & 0.470 & 0.517 & 15.58 \\

& AnySplat
& \underline{447K} & 0.167 & 0.819 & 24.07
& \underline{820K} & 0.148 & 0.842 & 25.10
& \underline{1142K} & 0.143 & 0.849 & 25.40 \\

& \textbf{\OURMODEL$_{\tau^+}$}
& \textbf{105K} & \underline{0.142} & \underline{0.847} & \underline{25.26}
& \textbf{210K} & \underline{0.130} & \underline{0.860} & \underline{25.75}
& \textbf{315K} & \underline{0.128} & \underline{0.862} & \underline{25.85} \\

& \textbf{\OURMODEL$_{\tau^-}$}
& \underline{447K} & \textbf{0.131} & \textbf{0.859} & \textbf{25.64}
& \underline{820K} & \textbf{0.120} & \textbf{0.869} & \textbf{26.10} 
& \underline{1142K} & \textbf{0.119} & \textbf{0.870} & \textbf{26.18} \\
\bottomrule
\end{tabular}
}
\vspace{-5pt}
\end{table*}

\begin{table*}[t]
\centering
\setlength{\tabcolsep}{0.09cm}
\renewcommand{\arraystretch}{0.95}
\caption{\textbf{Generalization to unseen datasets.} Models are trained on RE10K~\cite{zhou2018re10k} and evaluated on the unseen ACID dataset~\cite{liu2021acid}. We compare reconstruction quality under different numbers of input views (8, 16, and 24).}
\label{tab:crossdata_acid_nview}
\vspace{-6pt}

\resizebox{\textwidth}{!}{
\begin{tabular}{llcccc cccc cccc}
\toprule
& \multirow{2}{*}{\textbf{Method}}
& \multicolumn{4}{c}{\textbf{8 views}}
& \multicolumn{4}{c}{\textbf{16 views}}
& \multicolumn{4}{c}{\textbf{24 views}} \\
\cmidrule(lr){3-6} \cmidrule(lr){7-10} \cmidrule(lr){11-14}
& 
& \#GS$\downarrow$ & LPIPS$\downarrow$ & SSIM$\uparrow$ & PSNR$\uparrow$
& \#GS$\downarrow$ & LPIPS$\downarrow$ & SSIM$\uparrow$ & PSNR$\uparrow$
& \#GS$\downarrow$ & LPIPS$\downarrow$ & SSIM$\uparrow$ & PSNR$\uparrow$ \\
\midrule

\multirow{3}{*}{\shortstack[l]{\emph{Pose-Free}}}
& NoPoSplat
& 524K & 0.268 & 0.672 & 22.84
& 1049K & 0.294 & 0.644 & 22.14
& 1573K & 0.307 & 0.630 & 21.78 \\

& VicaSplat
& 524K & 0.271 & 0.656 & 22.57
& 1049K & 0.327 & 0.609 & 21.20
& 1573K & 0.356 & 0.581 & 20.37 \\

& SPFSplat
& 524K & 0.215 & 0.736 & 24.89
& 1049K & 0.214 & 0.743 & 25.07
& 1573K & 0.218 & 0.744 & 25.00 \\

\midrule

\multirow{4}{*}{\shortstack[l]{\emph{Uncalibrated}}}
& VicaSplat
& 524K & 0.315 & 0.588 & 21.31
& 1049K & 0.377 & 0.547 & 19.82
& 1573K & 0.399 & 0.529 & 19.22 \\

& AnySplat
& \underline{481K} & 0.248 & 0.696 & 23.30 
& \underline{906K} & 0.236 & 0.720 & 23.88
& \underline{1289K} & 0.234 & 0.727 & 24.04 \\

& \textbf{\OURMODEL$_{\tau^+}$}
& \textbf{52K} & \underline{0.239} & \underline{0.713} & \underline{24.28}
& \textbf{105K} & \underline{0.230} & \underline{0.726} & \underline{24.54}
& \textbf{315K} & \underline{0.216} & \underline{0.741} & \underline{24.72} \\

& \textbf{\OURMODEL$_{\tau^-}$}
& \underline{481K} & \textbf{0.204} & \textbf{0.744} & \textbf{24.83}
& \underline{906K} & \textbf{0.201} & \textbf{0.753} & \textbf{25.01}
& \underline{1289K} & \textbf{0.203} & \textbf{0.752} & \textbf{24.88} \\
\bottomrule
\end{tabular}
}
\vspace{-15pt}
\end{table*}
\subsection{Experimental Results}
\vspace{-3pt}
As shown in \cref{tab:re10k_multiviews} and \cref{tab:crossdata_acid_nview}, our method achieves competitive performance among uncalibrated baselines and remains competitive with pose-free and pose-required methods. When using a similar number of Gaussian primitives as the baselines, our approach achieves strong reconstruction performance. Notably, even when using only 10-28\% of the Gaussian primitives, our method still maintains competitive results. These results suggest that our explicit density control enables a compact yet high-fidelity 3D Gaussian representation, leading to improved efficiency without sacrificing reconstruction quality.

\cref{fig:qual} shows qualitative comparisons on RE10K under different input-view settings. Our method produces sharper structures and more faithful scene details compared to existing approaches. Notably, even with substantially fewer Gaussian primitives (24-29\%), our approach maintains high rendering quality while reducing blurring artifacts, resulting in visually more accurate renderings.

\cref{tab:nvs_2view_acid} demonstrates the novel view synthesis performance on the ACID dataset under the challenging two-view setting. Our method achieves superior performance compared to the uncalibrated baseline. Moreover, it remains competitive with both pose-required and pose-free approaches, including feed-forward 3DGS methods and other baselines, demonstrating robust performance even with only two sparse input views.
\begin{figure*}[t]
    \centering
    \includegraphics[width=0.9\textwidth]{Figures/qual.pdf}
    \vspace{-4pt}
    \caption{\textbf{Qualitative comparisons of novel view synthesis on RE10K~\cite{zhou2018re10k}} 
    \label{tab:qual}
    The bottom-right corner of each image shows the PSNR of the rendered view and the number of Gaussian primitives used for scene reconstruction. 
    Our method achieves high-quality rendering even with substantially fewer primitives, while consistently outperforming competing approaches.}
    \label{fig:qual}
    \vspace{-15pt}
\end{figure*}
\subsection{Ablation Studies and Analysis}
\vspace{-3pt}
\noindent{\textbf{Densification-Score-Guided Allocation.}}
We validate the effectiveness of our learned densification score for Gaussian allocation by comparing it with two alternative strategies, (a) and (b) in \cref{tab:ablation}.
For (a) \textit{random-based allocation}, we replace the learned densification score with a uniform score, which leads to random Gaussian allocation.
For (b) \textit{frequency-based allocation}, we compute a heuristic frequency score by resizing the input image to each scale level and applying a Sobel filter to emphasize high-frequency (edge) regions.
Both alternatives underperform, demonstrating that the learned densification score provides a more effective signal for allocating Gaussians than either random selection or simple frequency-based heuristics. 
Qualitatively, \cref{fig:densification_score_analysis} further illustrates the behavior of the learned densification score.
The predicted maps of the top example (a) assign higher scores to structurally and visually complex regions (e.g., object boundaries, textured areas, and fine details), indicating where increased Gaussian density is likely to be beneficial.
Moreover, as in the example below, overlapping areas observed from multiple context images tend to have lower scores, reflecting redundancy-awareness across views.
This encourages allocating Gaussians to regions that need more capacity while avoiding redundancy in well-covered areas, improving efficiency and fidelity under a fixed budget.

\noindent{\textbf{Level-Wise Gaussian Supervision.}}
Next, we examine the impact of level-wise Gaussian supervision during training. If we supervise only the final Gaussian set $\mathcal{G}_\tau$, the sampled threshold $\tau$ causes some scale-level Gaussians to be excluded from training in each iteration, resulting in less stable optimization. As shown in \cref{tab:ablation}, removing level-wise supervision clearly degrades performance (c), confirming the benefit of supervising multi-scale Gaussians during training.

\noindent{\textbf{Scene-Scale Regularization.}}
Finally, we ablate the scene-scale regularization loss in~\cref{tab:ablation} (d). Without this term, training becomes highly unstable, and the model fails to learn a meaningful reconstruction, leading to training failure. In the uncalibrated setting, where the model must jointly predict camera parameters and scene geometry, this simple regularizer is crucial for stable optimization.

\begin{table}[!t]
\centering
\setlength{\tabcolsep}{0.08cm}
\renewcommand{\arraystretch}{0.9}

%%%%%%%%%%%%%%%%%%%%%%%%%%%%%%%%%%%%%%%%%%%%%%%%%%%%
\begin{minipage}[t]{0.49\linewidth}
\caption{\textbf{Novel view synthesis performance comparison under 2-view setting.}
Results on ACID~\cite{liu2021acid} dataset.}
\label{tab:nvs_2view_acid}
\vspace{-5pt}

\resizebox{\linewidth}{!}{
% \resizebox{0.5\columnwidth}{!}{
\begin{tabular}{llcccc}
\toprule
& \multirow{2}{*}{\textbf{Method}}
& \multicolumn{4}{c}{\textbf{Average}} \\
\cmidrule(lr){3-6}
& 
& \#GS$\downarrow$ & LPIPS$\downarrow$ & SSIM$\uparrow$ & PSNR$\uparrow$ \\
% & Method & \#GS$\downarrow$ & LPIPS$\downarrow$ & SSIM$\uparrow$ & PSNR$\uparrow$ \\
\midrule

\multirow{4}{*}{\shortstack[l]{\emph{Pose-} \\ \emph{Required}}}
& pixelNeRF
& - & 0.533 & 0.561 & 20.323 \\

% & AttnRend
% & - & 0.287 & 0.730 & 24.475 \\

& pixelSplat 
& 131K & 0.195 & 0.779 & 25.819 \\

& MVSplat 
& 131K & 0.196 & 0.773 & 25.512 \\

% & DepthSplat
% & TODO & TODO & TODO & TODO
% & TODO & TODO & TODO & TODO
% & TODO & TODO & TODO & TODO
% & -- & -- & -- & -- \\
\midrule

\multirow{6}{*}{\shortstack[l]{\emph{Pose-Free}}}
& DUSt3R
& - & 0.447 & 0.411 & 16.286 \\

& MASt3R
& - &  0.461 & 0.409 &16.179 \\

& CoPoNeRF
& - & 0.406 & 0.606 & 20.950 \\

& NoPoSplat 
& 131K & 0.189 & 0.781 & 25.961 \\

& VicaSplat
& 131K & 0.201 & 0.757 & 25.439 \\

& SPFSplat
& 131K & 0.176 & 0.807 & 26.796 \\
% & FLARE & -- & -- & -- & -- \\
\midrule

\multirow{3}{*}{\shortstack[l]{\emph{Uncalibrated}}}
& VicaSplat 
& 131K & 0.218 & 0.726 & 24.548 \\
% & AnySplat$^\dagger$ & 124K & 21.14 & 0.735 & 0.235 \\

& \textbf{\OURMODEL}$_{\tau^+}$
& \textbf{52K} & 0.188 & 0.784 & 26.028 \\

& \textbf{\OURMODEL}$_{\tau^-}$
& 131K & \textbf{0.176} & \textbf{0.794} & \textbf{26.282} \\
\bottomrule
\end{tabular}
}
\end{minipage}
\hspace{0.02\linewidth}
\begin{minipage}[t]{0.47\linewidth}
    \captionsetup{type=table}
    \caption{\textbf{Ablation Studies.} Experiments are conducted with 24 input views, where the Gaussian budget is fixed to 20\% of the maximum number of Gaussians. (a)-(b) compare different Gaussian allocation strategies, (c) removes the level-wise Gaussian supervision during training, and (d) removes the scene-scale regularization term.}
    \label{tab:ablation}
    \vspace{-5pt}
    \centering
    \setlength{\tabcolsep}{0.12cm}
    \renewcommand{\arraystretch}{0.95}
    \resizebox{\linewidth}{!}{
    \begin{tabular}{clccc}
        \toprule
        & \textbf{Variant} & \textbf{LPIPS}$\downarrow$ & \textbf{SSIM}$\uparrow$ & \textbf{PSNR}$\uparrow$ \\
        \midrule
        (a) & Rand. based allocation            & 0.194 & 0.828 & 24.68 \\
        (b) & Freq. based allocation   & 0.160 & 0.841 & 25.36 \\
        (c) & w/o level-wise GS train            & 0.192 & 0.813 & 24.25 \\
        (d) & w/o scene scale reg.        & 0.712 & 0.006 & \phantom{0}4.82 \\
        (e) & \textbf{Ours}               & \textbf{0.143} & \textbf{0.854} & \textbf{25.47} \\
        \bottomrule
    \end{tabular}
    }
\end{minipage}
\vspace{-15pt}
\end{table}

